\section{Taylor's Formula}
\index{formula!di Taylor}

\begin{definition}[Taylor polynomial]
\mymark{***}
Let $A\subset \RR$, $f\colon A \to \RR$ be a function
and let $x_0\in A$ be a point where the $n$-th derivative
(and hence also the preceding derivatives) is defined.

Then the polynomial:
\[
  P(x) = \sum_{k=0}^n \frac{f^{(k)}(x_0)}{k!}(x-x_0)^k
\]
is called the \emph{Taylor polynomial}%
\mymargin{Taylor polynomial}%
\index{polynomial!Taylor}
of order $n$
of the function $f$,
centered at $x_0$.
\end{definition}

In the particular case where $x_0=0$, the Taylor polynomial is also called the
\emph{MacLaurin polynomial}%
\mymargin{MacLaurin polynomial}%
\index{polynomial!MacLaurin}.

\begin{example}
  The Taylor polynomial of order $3$ for the function $f(x) = \sin x$ 
  centered at $x_0=0$ is:
  \[
    P(x) = x - \frac{x^3}{6}.
  \]
  Indeed, we have $f(0) = \sin(0) = 0$, $f'(0) = \cos(0) = 1$, 
  $f''(0) = -\sin(0) = 0$, and $f'''(0) = -\cos(0) = -1$.
  Clearly, $1!=1$ and $3!=6$. 
\end{example}

\begin{theorem}[characterization of the Taylor polynomial]%
\label{th:caratterizzazioneTaylor}%
\mymark{***}%
The Taylor polynomial of a function $f$, of order $n$, centered at $x_0$ is the
unique polynomial $P$ of degree not exceeding $n$ such that
\[
  P^{(k)}(x_0) = f^{(k)}(x_0), \qquad \forall k \in \ENCLOSE{0,\dots,n}.
\]
\end{theorem}
%
\begin{proof}
  Every polynomial of degree not exceeding $n$ can be written in the form:
\[
  P(x) = \sum_{k=0}^{n} a_k (x-x_0)^k
\]
(it suffices to note that if $P(x)$ is a polynomial, then $Q(t) = P(x_0+t)$ is
also a polynomial)
and its derivatives are:
\begin{align*}
   P^{(j)}(x)
   = \sum_{k=j}^n a_k \cdot \frac{k!}{(k-j)!}\cdot (x-x_0)^{k-j}.
\end{align*}
For $x=x_0$, the only non-zero term is the one with $k=j$, thus
\[
  P^{(k)}(x_0) = k! \cdot a_k.
\]
Therefore, $P^{(k)}(x_0) = f^{(k)}(x_0)$ if and only if $a_k = f^{(k)}(x_0)/k!$,
which means $P$ is the Taylor polynomial of $f$.
\end{proof}

\begin{remark}[Taylor polynomial of the derivative]
  \label{rem:taylor_derivata}%
If $P_n$ is the Taylor polynomial centered at $x_0$ of order $n$ for $f$, then
we have
\begin{align*}
P_n'(x) &= \sum_{k=1}^{n} \frac{f^{(k)}(x_0)}{k!}\cdot k(x-x_0)^{k-1} \\
  &= \sum_{k=1}^{n} \frac{f^{(k)}(x_0)}{(k-1)!} (x-x_0)^{k-1} \\
  &= \sum_{j=0}^{n-1} \frac{f^{(j+1)}(x_0)}{j!} (x-x_0)^j.
\end{align*}
Thus, it is verified that $P'_n$ is the Taylor polynomial of order $n-1$ for
$f'$.
In short: the Taylor polynomial of order $n-1$ of the derivative is the
derivative
of the Taylor polynomial of order $n$ of the function.

Conversely, if
\[
   \sum_{k=0}^n a_k (x-x_0)^k
\]
is the Taylor polynomial of order $n$ of the derivative $f'(x)$, then
the Taylor polynomial of order $n+1$ of $f$ will be%
\mynote{This is an \emph{antiderivative} (or \emph{primitive}) of the given polynomial, 
see definition~\ref{def:primitiva}.}
\begin{equation}\label{eq:4993725}
  P_{n+1}(x) = f(x_0) + \sum_{k=0}^{n} \frac{a_{k}}{k+1}\cdot (x-x_0)^{k+1}
\end{equation}
since
\[
  \frac{f^{(k+1)}(x_0)}{(k+1)!}
  =\frac{(f')^{(k)}(x_0)}{(k+1)!}
  = \frac{a_k \cdot k!}{ (k+1)!}
  = \frac{a_k}{k+1}.
\]
\end{remark}

\begin{theorem}[Taylor's formula with Lagrange remainder]
\label{th:taylor_lagrange}%
\mymark{**}%
\index{formula!di Taylor!with Lagrange remainder}%
\index{theorem!of Taylor with Lagrange remainder}%
\index{Taylor!Lagrange remainder}
\mymargin{Taylor with Lagrange remainder}
Let $I\subset \RR$ be an interval, $x_0\in I$
and $f\in C^{n+1}(I)$.
% \mynote{It is understood that $I$ cannot be empty, because $x_0\in I$, 
% and $I$ cannot be a singleton because if $I=\ENCLOSE{x_0}$
% no function can be differentiable at $x_0$ 
% since $x_0$ is not an accumulation point for $I$.
% }%
Let $P$ be the Taylor polynomial of $f$ of order $n$ centered at $x_0$.
For every $x\in I\setminus\ENCLOSE{x_0}$ there exists $y\in (x_0,x)$
\mynote{%
in this proof, intervals are considered 
non-oriented: $[a,b] = [b,a]$, $(a,b)=(b,a)$ if $a>b$.}
such that
\[
  f(x) = P(x) + \frac{f^{(n+1)}(y)}{(n+1)!}(x-x_0)^{n+1}.
\]
\end{theorem}
%
\begin{proof}\mymark{*}%
The thesis can be written in the form 
\[
  \frac{f(x) - P(x)}{(x-x_0)^{n+1}} 
  = \frac{f^{(n+1)}(y)}{(n+1)!}.
\]

We proceed with an induction proof on $n$. 
For $n=0$, we have that $f$ is differentiable (indeed it is $C^1$, but this is
an
excessive hypothesis) and the thesis is
\begin{equation}\label{eq:655643}
  \frac{f(x) - f(x_0)}{x-x_0}  = f'(y)  
\end{equation}
which coincides with Lagrange's theorem~\ref{th:lagrange}.

Suppose then by induction 
that the theorem is true with $n-1$ instead of 
$n$ and consider $f\in C^{n+1}$ with Taylor polynomial $P$
of order $n$ centered at $x_0$. 
The function $f'$ is in $C^{n}$ and its Taylor polynomial 
of order $n-1$ is $P'$ (observation~\ref{rem:taylor_derivata}).
To obtain~\eqref{eq:655643}, observe now that, since $P(x_0)=f(x_0)$:
\[
  \frac{f(x)-P(x)}{(x-x_0)^{n+1}}  
  = \frac{f(x) - P(x) - (f(x_0)-P(x_0))}{(x-x_0)^{n+1} - (x_0-x_0)^{n+1}}
\]
and applying Cauchy's theorem~\ref{th:cauchy}
we find that there exists $t\in (x_0,x)$ such that
\[
  \frac{f(x)-P(x)}{(x-x_0)^{n+1}}  
  = \frac{f'(t)-P'(t)}{(n+1)(t-x_0)^n}. 
\]
Now we can apply the inductive hypothesis to the function 
$f'$ with Taylor polynomial $P'$
and putting $t$ in place of $x$ to obtain 
the existence of $y\in (x_0,t)\subset(x_0,x)$ such that 
\[
  \frac{f'(t)-P'(t)}{(n+1)(t-x_0)^n}
  = \frac{(f')^{(n)}(y)}{(n+1)\cdot n!}.
\]
Observing that $(f')^{(n)} = f^{(n+1)}$ and that $(n+1)\cdot n!=(n+1)!$ 
we obtain the thesis.
\end{proof}

\begin{example}[sum of the alternating harmonic series]
We have:
  \[
    \sum_{k=1}^{+\infty} \frac{(-1)^{k-1}}{k} = \ln 2.
  \]
\end{example}
\begin{proof}
Consider the function $f(x) = \ln x$.
Observe that
\begin{gather*}
  f'(x) = x^{-1},\qquad
  f''(x) = -x^{-2},\qquad
  f'''(x) = 2x^{-3}, \\
  f^{(4)}(x) = -3\cdot 2 x^{-4}, \qquad
  \dots. \qquad
  f^{(k)}(x) = \frac{(-1)^{k-1}(k-1)!}{x^k}
\end{gather*}
and thus, for every $k\ge 1$,
\[
  \frac{f^{(k)}(1)}{k!} = \frac{(-1)^{k-1}}{k}
\]
from which it follows that the Taylor polynomial of order $n$
centered at $x_0=1$ is
\[
P(x) = \sum_{k=1}^n (-1)^{k-1}\frac{(x-1)^k}{k}.
\]
Taylor's formula with Lagrange remainder tells us that for
every $x>1$ there exists $y\in(1,x)$ such that
\[
  \ln(x) - \sum_{k=1}^n (-1)^{k-1}\frac{(x-1)^k}{k} 
  = \frac{(-1)^n(x-1)^{n+1}}{(n+1)y^{n+1}}.
\]
If $x=2$, we have $y>1$ and $(x-1)^{n+1}=1$, thus
the right-hand side tends to zero as $n\to +\infty$ 
(note that $y$ also depends on $n$). 
But for $x=2$ and $k\to +\infty$, the left-hand side tends to 
\[
 \ln 2 - \sum_{k=1}^\infty \frac{(-1)^{k-1}}{k}
\]
and thus we obtain the desired result.
\end{proof}

In the following theorem, we will use a new notation. 
\begin{definition}[little-$o$ notation]
We will write
\mymargin{little-$o$}%
\index{little-$o$}%
\[
  f(x) = P(x) + o(g(x)), \qquad \text{as $x\to x_0$}
\]
(read: ``$f(x)$ is equal to $P(x)$ plus a little-$o$ of $g(x)$'')
if it holds that
\[
  f(x) - P(x) \ll g(x)
  \qquad\text{or equivalently}\qquad
  \lim_{x\to x_0} \frac{f(x)-P(x)}{g(x)} = 0.
\]
\end{definition}

The quantity $o(g(x))$ represents an unknown function
that we know to be negligible (in the sense just described)
compared to the function $g(x)$.
This notation is very convenient but must be used with caution:
in chapter~\ref{sec:landau}, we will see how to rigorously formalize 
this notation.

\begin{theorem}[Taylor's formula with Peano remainder]
\label{th:taylor_peano}%
\mymark{***}%
\index{formula!di Taylor!with Peano remainder}%
\index{theorem!of Taylor with Peano remainder}%
\index{Taylor!Peano remainder}%
Let $I\subset \RR$ be an interval, $x_0\in I$, 
and suppose that $f$ is differentiable $n-1$ times on all of $I$
and that $f^{(n)}(x_0)$ exists.
\mynote{In particular, the hypotheses are satisfied if $f\in C^n(I)$.}%
Let $P$ be the Taylor polynomial of $f$ of order $n>0$ centered at $x_0$. 
Then we have
\begin{equation}\label{eq:taylor_peano}
  \lim_{x\to x_0}\frac{f(x) - P(x)}{(x-x_0)^n} = 0
\end{equation}
or, expressed more informatively:
\begin{equation}\label{eq:Taylor}
  f(x) = P(x) + o((x-x_0)^n).
\end{equation}

Conversely, if $P(x)$ is any polynomial of degree less than or equal to $n$, and
the
formula~\eqref{eq:Taylor} holds, then $P$ is the Taylor polynomial of $f$.
\end{theorem}
%
\begin{proof}
\mymark{**}%
\mynote{Here we also use the convention $[a,b] = [b,a]$ when $b<a$.}
%
We proceed by induction in a manner similar to the proof 
of theorem~\ref{th:taylor_lagrange}.
%
For $n=1$, we have $P(x)=f(x_0) + f'(x_0)\cdot (x-x_0)$ and thus 
\[
 \frac{f(x)-P(x)}{x-x_0} = \frac{f(x) - f(x_0)}{x-x_0} - f'(x_0) 
 \to f'(x_0) - f'(x_0) = 0  \qquad\text{as $x\to x_0$.}
\]

Suppose now that $n>1$ and that the theorem is true
when $n-1$ is substituted for $n$. We have 
\[
  \frac{f(x)-P(x)}{(x-x_0)^n} 
  = \frac{f(x)-P(x)-(f(x_0)-P(x_0))}{(x-x_0)^n - (x_0-x_0)^n}
\]
and, since $P(x_0) = f(x_0)$, we can apply Cauchy's theorem~\ref{th:cauchy}
to the functions $f(x)-P(x)$ and $(x-x_0)^n$,
to obtain that there exists $t\in (x_0,x)$ such that
\[
  \frac{f(x)-P(x)}{(x-x_0)^n} = \frac{f'(t)-P'(t)}{n\cdot (t-x_0)^{n-1}}.
\]
We now apply the inductive hypothesis to the function $f'(t)$ whose 
Taylor polynomial of order $n-1$ is precisely $P'(t)$. We thus obtain 
that 
\[
 \lim_{t\to x_0} \frac{f'(t)-P'(t)}{(t-x_0)^{n-1}} = 0 
\]
and since for $x\to x_0$, being $t\in(x_0,x)$, we also have 
$t\to x_0$, we obtain, as we wanted to demonstrate,
\[
 \lim_{x\to x_0} \frac{f(x)-P(x)}{(x-x_0)^n} 
 = \lim_{t\to x_0} \frac{f'(t)-P'(t)}{n\cdot (t-x_0)^{n-1}} = 0. 
\]

\emph{Reverse implication.}
To complete the proof, we must show that the Taylor polynomial $P$ is the unique
polynomial of degree not exceeding $n$ that satisfies
condition~\eqref{eq:taylor_peano}.
If there were another polynomial $Q$ with the same property, the difference 
$R=(f-P)-(f-Q)=Q-P$ would have to be a polynomial of degree 
not exceeding $n$ that satisfies the condition 
\begin{equation}\label{eq:496390}
  \lim_{x\to x_0}\frac{R(x)}{(x-x_0)^n}=0.  
\end{equation}
Setting 
\[
  R(x) = \sum_{k=0}^n a_k (x-x_0)^k  
\]
we observe that $R(x)\to R(x_0) = a_0$ as $x\to x_0$. 
Then it must necessarily be $a_0 = 0$, 
otherwise the limit~\eqref{eq:496390} would be infinite.
But if $a_0=0$, we can factor out $x-x_0$ in both the numerator 
and denominator and from~\eqref{eq:496390} obtain 
\[
  \lim_{x\to x_0}\frac{\displaystyle \sum_{k=1}^n a_k (x-x_0)^{k-1}}{(x-x_0)^{n-1}}=0.  
\]
Iterating the process, we discover that all the coefficients $a_k$ 
must be zero, and thus $Q=P$.
\end{proof}

\begin{exercise}
Calculate:
\[
  \lim_{x\to 0} \frac{1-\cos(x^2)}{x^4}.
\]
\end{exercise}
\begin{proof}[Solution.]
The Taylor polynomial of order $2$ for the function $f(x) = \cos x$ 
centered at $x_0=0$ is 
\[ 
  P(x)
   = \cos 0 - \sin 0\cdot x -\frac{\cos 0}{2} \cdot x^2 
   = 1 - \frac{x^2}{2}.
\]
Theorem~\ref{th:taylor_peano} (Taylor's formula with Peano remainder)
tells us that we have  
\begin{equation}\label{eq:498101}
   \lim_{x\to 0} \frac{\cos x - 1 + \frac {x^2}{2}}{x^2} = 0
\end{equation}
or, using the little-$o$ notation:
\begin{equation}\label{eq:498102}
  \cos x = 1 - \frac{x^2}{2} + o(x^2)
  \qquad \text{as $x\to 0.$}
\end{equation}
Substituting $x^2$ in place of $x$ in the limit~\eqref{eq:498101} 
we consequently obtain:
\begin{equation}\label{eq:498103}
  \lim_{x\to 0}\frac{\cos x^2 -1 + \frac{x^4}{2}}{x^4} = 0
\end{equation}
which can be expressed more informatively as:
\[
  \cos (x^2) = 1 - \frac{x^4}{2} + o(x^4).
\]
This last relation can be obtained directly 
from~\eqref{eq:498102} by substituting $x^2$ for $x$, 
as we will see in chapter~\ref{sec:landau}.
Then as $x\to 0$, we have:
\[
  \frac{1-\cos(x^2)}{x^4} 
  = \frac{\frac{x^4}{2} - o(x^4)}{x^4} 
  = \frac 1 2 - \frac{o(x^4)}{x^4} \to \frac 1 2
\]
since, by definition of little-$o$, we have
\[
  \lim_{x\to 0 } \frac{o(x^4)}{x^4} = 0.
\]
\end{proof}

\begin{example}
Let $f(x) = \cos(x^2)$. Calculate $f''''(0)$.
\end{example}
\begin{proof}[Solution.]
We have seen in the previous example that as $x\to 0$, we have 
\[
  \cos(x^2) = 1 - \frac{x^4}{2} + o(x^4).
\]
Thanks to the second part of theorem~\ref{th:taylor_peano} (Taylor's formula
with Peano remainder)
we know that the polynomial 
\[
  P(x) = 1 - \frac{x^4}{2}
\]
is the Taylor polynomial of order $4$ centered at $x_0=0$ for the function
$f(x)$.
On the other hand, by the definition of the Taylor polynomial, we have
\[
  P(x) = f(0) + f'(0) \cdot x + \frac{f''(0)}{2} \cdot x^2 
  + \frac{f'''(0)}{6} \cdot x^3 + \frac{f''''(0)}{24} \cdot x^4
\]
and therefore, by the principle of identity of polynomials, we deduce
\[
  f(0) = 1,\qquad
  f'(0) = 0,\qquad
  f''(0) = 0,\qquad
  f'''(0) = 0,\qquad
  f''''(0) = -12.
\]
\end{proof}

\begin{definition}[real binomial coefficient]
\label{def:binomiale_reale}%
\mymark{***}%
Given $\alpha \in \RR$, $k \in \NN$, we define
\[
 {\alpha \choose k } 
 = \frac{\alpha \cdot (\alpha-1) \cdots (\alpha -k +1)}{k!}
 = \prod_{j=1}^k \frac{\alpha-j+1}{j}.
\]
We observe that if $\alpha \in \NN$, this definition coincides
with the definition given in chapter~\ref{ch:binomiale}.
\end{definition}

\begin{theorem}[Taylor expansions of some elementary functions]
\label{th:sviluppi_taylor}%
\mymark{***}%
For every $n\in \NN$, we have, as $x\to 0$,
the relations reported in table~\ref{tb:taylor}
on page~\pageref{tb:taylor}.
\end{theorem}
\begin{table}
\begin{align*}
e^x &= \sum_{k=0}^n \frac{x^k}{k!} + o(x^n) \\
  &= 1 + x + \frac{x^2}{2} + \frac {x^3}{6} + \dots + \frac{x^n}{n!} + o(x^n) \\
\sin x &= \sum_{k=0}^n (-1)^k\frac{x^{2k+1}}{(2k+1)!} + o(x^{2n+2}) \\
  &= x - \frac{x^3}{6} + \dots + (-1)^{n} \frac{x^{2n+1}}{(2n+1)!}  +
 o(x^{2n+2})\\
 \cos x &= \sum_{k=0}^n (-1)^k \frac{x^{2k}}{(2k)!} + o(x^{2n+1}) \\
      &= 1 - \frac{x^2}{2} + \frac{x^4}{24} - \dots + (-1)^n\frac{x^{2n}}{(2n)!}
   + o(x^{2n+1})\\
 (1+x)^\alpha &= \sum_{k=0}^n {\alpha \choose k} x^k + o(x^n)\\
        &= 1 + \alpha x + \frac{\alpha (\alpha-1)}{2} x^2 + \dots + {\alpha
    \choose n} x^n + o(x^n) \\
  \ln (1+x) &= \sum_{k=1}^n (-1)^{k-1} \frac{x^k}{k} + o(x^n) \\
                  &= x - \frac{x^2}{2} + \frac{x^3}{3} - \dots +
         (-1)^{n-1}\frac{x^n}{n} + o(x^n) \\
  \arctg x &= \sum_{k=0}^n (-1)^k \frac{x^{2k+1}}{2k+1} + o(x^{2n+1})\\
        &= x - \frac{x^3}{3} + \frac{x^5}{5} - \dots + (-1)^n \frac{x^{2n+1}}{2n
    +1} + o(x^{2n+1})\\
  \arcsin x &= x + \frac{x^3}{6} + \frac{3}{40} x^5 + \dots + \frac{(2n-1)!!}{(2n)!!} \frac{x^{2n+1}}{2n+1} + o(x^{2n+1}) \\
  \arccos x &= \frac \pi 2 - \arcsin x\\
  \tg x &= x + \frac{x^3}{3} + \frac{2}{15} x^5 + o(x^5).
\end{align*}
\caption{Taylor expansions, as $x\to 0$, of some elementary functions.
\index{Taylor!polynomials of elementary functions}%
\index{expansion!Taylor polynomial}%
See theorem~\ref{th:sviluppi_taylor}.}
\label{tb:taylor}%
\end{table}
%
\begin{proof}
\mymark{**}
By definition, the coefficient of the term $x^k$
in the Taylor polynomial of $f(x)$ is nothing but $a_k=f^{(k)}(0)/k!$
If $f(x) = e^x$, then $f^{(k)}(x) = e^x$ and thus $f^{(k)}(0) = 1$. We find the
coefficients $a_k = 1/k!$.

If $f(x) = \sin x$, we have $f'(x)=\cos x$, $f''(x) =-\sin x$, $f'''(x) = -\cos
x$, and $f^{(4)}(x) = \sin x$... and then the derivatives repeat every four
iterations. Evaluating the derivatives at $x=0$, we obtain the sequence $0, 1,
0, -1, \dots$ which repeats indefinitely. We thus obtain the indicated
expansion. A similar argument can be made for $f(x) = \cos x$.

If $f(x) = (1+x)^\alpha$, we obtain $f'(x) = \alpha (1+x)^{\alpha -1}$, $f''(x)
= \alpha (\alpha -1) (1+x)^{\alpha -2}$, and so on...
Evaluating the derivatives at $x=0$, we obtain the sequence: $1$, $\alpha$,
$\alpha(\alpha-1)$, $\alpha (\alpha-1)(\alpha -2)$\dots from which, dividing by
$k!$, we obtain that the coefficients of the Taylor polynomial are the binomial
coefficients ${\alpha \choose k}$.

If $f(x) = \ln (1+x)$, we observe that $f(0)=0$, then we have $f'(x) =
(1+x)^{-1}$, and the subsequent derivatives coincide with the derivatives of
$(1+x)^\alpha$ with $\alpha=-1$: the coefficients (except the first which is
zero) coincide with the binomial coefficients
\[
{-1 \choose k} = \frac{(-1)(-2)(-3) \dots (-k)}{k!} = (-1)^k.
\]

Similarly, if $f(x) = \arctg x$, we observe that $f(0) = 0$ and $f'(x) =
(1+x^2)^{-1}$. From what we have already seen, we know that
\[
 (1+y)^{-1} = \sum_{k=0}^n (-1)^k y^k + o(y^n)
\]
from which, substituting $y=x^2$, we have
\[
 f'(x) = (1+x^2)^{-1} = \sum_{k=0}^n (-1)^k x^{2k} + o(x^{2n}).
\]
Using the second part of Theorem~\ref{th:taylor_peano} (Taylor's formula with
Peano remainder),
we have verified that equation \eqref{eq:Taylor} holds
for the polynomial $P(x) = \sum_{k=0}^n (-1)^k x^{2k}$.
Thus $P(x)$ is the Taylor polynomial of degree $2n$ of $f'(x)$.
But the Taylor polynomial of the derivative is the derivative of the Taylor
polynomial, so we can use \eqref{eq:4993725} to obtain
the polynomial reported in the table.

A similar method is used for $f(x)=\arcsin x$, observing that
\begin{align*}
  f'(x) &= \frac{1}{\sqrt{1-x^2}}
  = (1-x^2)^{-\frac 1 2} \\
  &= \sum_{k=0}^n {-\frac 1 2 \choose k} (-x^2)^k + o(x^{2n}) \\
    &= \sum_{k=0}^n \frac{\enclose{-\frac 1 2}\cdot \enclose{-\frac 1 2 -
  1}\cdots \enclose{-\frac 1 2 - k+1}}{k!}(-1)^k x^{2k} + o(x^{2n}) \\
    &= \sum_{k=0}^n \frac{\frac{-1}{2} \cdot \frac{-3}{2}\cdots
  \frac{-(2k-1)}{2}}{k!}(-1)^k x^{2k} + o(x^{2n}) \\
  &= \sum_{k=0}^n \frac{(2k-1)!!}{2^k \cdot k!}x^{2k} + o(x^{2n}) \\
  &= \sum_{k=0}^n \frac{(2k-1)!!}{(2k)!!}x^{2k} + o(x^{2n})
 \end{align*}
we have found the Taylor polynomial of $f'(x)$ and using \eqref{eq:4993725}
we obtain the formula reported in the table.

Regarding the function $f(x) = \tg x$, we limit ourselves to explicitly
calculating the first terms:
\begin{align*}
  f(x) &= \tg x
  & f(0)&=0\\
  f^{(1)} &= 1+ f^2
  & f^{(1)}(0) &= 1\\
  f^{(2)} &= 2ff'
  & f^{(2)}(0) &= 0 \\
  f^{(3)} &= 2(f')^2 + 2f f''
  & f^{(3)}(0) &= 2 \\
  f^{(4)} &= 4 f' f'' + 2 f' f'' + 2 f f'''
   & f^{(4)}(0) &= 0 \\
   f^{(5)} &= 6 (f'')^2 + 6 f' f''' + 2 f' f''' + 2 f f^{(4)}&&\\
           &= 6 (f'')^2 + 8 f' f''' + 2 f f^{(4)}
    & f^{(5)}(0) &= 16.
\end{align*}
We thus obtain the coefficients:
  $a_0 = 0$, $a_1 = 1$, $a_2 = 0$, $a_3 = 2/3! = 1/3$, $a_4=0$, $a_5 = 16/ 5! =
 2/15$.
\end{proof}